\documentclass[leqno, 12pt]{jsarticle}
\usepackage{amssymb}
\usepackage{amsthm}
\usepackage{amsmath}
\usepackage{comment}
	%可換図式用
		%\usepackage[all]{xy}
	%重くなるので使わいときはコメント化しておく。
%定理環境 thmはあまり使わずpropをなるべく使う。
%主張は基本的に証明の中で使い、そのため番号を付けない。
%注意環境を付けてみた。
\newtheorem{thm}{定理}
\newtheorem{prop}[thm]{命題}
\newtheorem{cor}[thm]{系}
\newtheorem{lem}[thm]{補題}
\newtheorem{df}[thm]{定義}
\newtheorem{exam}[thm]{例}
\newtheorem{fact}[thm]{事実}
\newtheorem*{claim}{主張}
\newtheorem*{rmk}{注意}
\renewcommand{\proofname}{{\bf 証明}}
%矢印たち。
%\toは\rightarrowと同じ。\getsは\leftarrowと同じ。同値は\iff
%上向き矢はuparrow逆はdownarrow
\newcommand{\To}{\Longrightarrow}
\newcommand{\Gets}{\Longleftarrow}
%黒板太字
\newcommand{\nn}{\mathbb{N}}
\newcommand{\zz}{\mathbb{Z}}
\newcommand{\qq}{\mathbb{Q}}
\newcommand{\rr}{\mathbb{R}}
\newcommand{\cc}{\mathbb{C}}
\newcommand{\ff}{\mathbb{F}}
%無限大
\newcommand{\mgn}{\infty}
%ギリシャン
\newcommand{\dpst}{\displaystyle}
\newcommand{\ep}{\varepsilon}
\newcommand{\al}{\alpha}
\newcommand{\be}{\beta}
\newcommand{\de}{\delta}
\newcommand{\la}{\lambda}
\newcommand{\La}{\Lambda}
\newcommand{\ga}{\gamma}
\newcommand{\lil}{\la\in \La}
%商集合
\newcommand{\qt}[2]{#1/\! #2}
\newcommand{\leg}[2]{\left(\frac{#1}{#2}\right)}
%enumerate環境
\renewcommand{\labelenumi}{{\rm (\arabic{enumi})}}
%以降alignじゃなくてenumerate を使え。
%なんか演算
\DeclareMathOperator{\card}{card}
\DeclareMathOperator{\supp}{supp}
%なんか演算２
\DeclareMathOperator{\bd}{Bdry}
\DeclareMathOperator{\cl}{CL}
\DeclareMathOperator{\nai}{Int}
\DeclareMathOperator{\di}{diam}

%差集合は\setminusで統一する。等号の否定は\neq
%真部分集合は\subsetneqq　偏微分は\partial
%ディスプレイスタイル。\dfracでディスプレイ分数
%空集合は\Oを使う！！！！！！！！！！！！

\title{平方剰余の相互法則}
\author{ナンブキトラ}
\date{\today}
			\begin{document}
\maketitle

この文章では
平方剰余の補充則と
相互法則の可能な限り代数的な証明を
紹介する．


まず最初に
\TeX で
\[
a\equiv b \mod p
\]
という式を入力するのがとても面倒だと思ったので，
この文章では，この式を
\[
a=_{p} b
\]
と書くことにする．
$x\in \zz$と素数$p$に対して
$x=_pa^2$となる
$a$が存在する時$x$を$p$の平方剰余という．
平方剰余でない数を平方非剰余と呼ぶ．

可換環$R$に対して$R^{\times}$で，
$R$の可逆元全体からなる群を表す．

素数$p$に対して
$\ff_p$で$\zz/p\zz$を表す．
商写像
$\zz\to \ff_p$を
$\pi_p$で表す．
$\ff_p$は標数が$p$の体のになる．
また$\ff_p^{\times}$は巡回群になる．
この巡回群の生成元を$p$の原始根という．

写像
$s:\ff_p^{\times}\to \ff_p^{\times}$を
$s(x)=x^2$で定義する．
$\ff_p^{\times}$は可換なので，
$f$は準同型になる．
よって特に$s(\ff_p^{\times})$
は$\ff_p^{\times}$の部分群になる．
この群を$S_p$と書くことにする．
乗法群
$\ff_p^{\times}$は可換なので
もちろん
$S$は正規部分群である．
$p\neq 2$
ならば，
体の一般論から
$a\in \ff_p$に対して
$x^2=a$
となる
$x\in \ff_p$
はちょうど2つあるので，
$S$
の位数は
$(p-1)/2$である．
よって
$p\neq 2$のとき，
剰余群
$\ff_p^{\times}/S_p$
は位数が2の
群になる．
位数が$2$の群は全て同型なので，
同型写像
$e_p:\ff_p^{\times}/S_p\to \{1, -1\}$
をとる．
このとき
$a\in \zz$
に対してルジャンドル記号を
$(a|p):=e_p\circ s\circ \pi_p(a)$
と定義する．
$\dpst \left (\frac{a}{p}\right)$
は
$(a|p)$
と同じ意味であると約束する．
ルジャンドル記号は次のように明示できる．
\[
\left(\frac{a}{p}\right)=
\begin{cases}
1 & \text{$a$は$p$の平方剰余，}\\
-1 & \text{$a$は$p$の平方非剰余．}
\end{cases}
\]

つまり平方剰余が
$1$であるか$-1$であるかに応じて平方剰余かそうでないかがわかる．


定義から次のことがわかる．
\begin{prop}
任意の$a\in \zz$に対して，
\[
\left(\frac{a}{p}\right)=\left(\frac{a+p}{p}\right)=\left(\frac{a-p}{p}\right)
\]
が成り立つ．
\end{prop}


今からルジャンドル記号の計算方法を説明してゆく．
まず次のオイラーの基準が基本的である．
\begin{prop}
$p$を奇素数とする．
このとき任意の
$a\in \zz$に対して
\[
\leg{a}{p}=_p a^{\frac{p-1}{2}}
\]
が成り立つ．
\end{prop}
\begin{proof}
最初に注意するが，
$p$が奇数なので
$\frac{p-1}{2}$は整数である．
まずフェルマーの小定理から任意の
$x\in \ff_{p}^{\times}$に対して
\[
x^{p-1}=1
\]
が成り立つ．
つまり任意の
$a\in \ff_p^{\times}$
は多項式
$x^{p-1}-1$
の根になっている．
さらに
\begin{equation}\label{eq:siki}
x^{p-1}-1=\left(x^{\frac{p-1}{2}}-1\right)\left(x^{\frac{p-1}{2}}+1\right)
\end{equation}
と因数分解できる．
さて$S_p=\{y^2\mid y\in \ff_p^{\times}\}$
でに注意すると
任意の$x\in S_p$に対して
$x=y^2$となる$y\in \ff_p$が存在する．
再びフェルマーの小定理から
$x^{\frac{p-1}{2}}=y^{p-1}=1$なので，
$x^{\frac{p-1}{2}}-1=0$である．
体の一般論から多項式
$x^{\frac{p-1}{2}}-1$の根は高々$\frac{p-1}{2}$個であり
$\card(S_p)=\frac{p-1}{2}$なので
\begin{equation}\label{eq:siki2}
S_p=\{x\mid x^{\frac{p-1}{2}}-1=0\}
\end{equation}
である.
任意の$x\in \ff_p\setminus S_p$は平方非剰余の剰余類だが，
$x^{p-1}-1=0$を満たすということと
式(\ref{eq:siki}), (\ref{eq:siki2})から，
$x^{\frac{p-1}{2}}+1=0$
を満たさなければならない．（体には零因子は$0$しかない．）
よってオイラーの基準が示された．
\end{proof}


このオイラーの基準を使って直ちに平方剰余の相互法則の第一補充則がわかる．
\begin{prop}
奇素数$p$と整数について，
\[
\leg{-1}{p}=(-1)^{\frac{p-1}{2}}
\]
が成り立つ．
\end{prop}
\begin{proof}
オイラーの基準から，
$(a|p)=_p(-1)^{\frac{p-1}{2}}$
である．
この両辺は
$1$か
$-1$の値しか取らず，
また$p$が奇素数であることから
$2\neq_p0$なので実際の等式として
$(a|p)=(-1)^{\frac{p-1}{2}}$を得る．
\end{proof}

平方剰余の相互法則の第二補充則は次のものである．
\begin{prop}
奇素数$p$に対して
\[
\leg{2}{p}=(-1)^{\frac{p^2-1}{8}}
\]
\end{prop}
\begin{proof}
$A= \{1, \dots, \frac{p-1}{2}\}$とおく．
任意の$a\in A$についてもちろん
\[
a=((-1)^{a}a)(-1)^a
\]
が成り立つ．
よって
\begin{align*}
\prod_{a\in A}a&=\prod_{a\in A}((-1)^aa)(-1)^a\\
&=\prod_{a\in A}((-1)^aa)\prod_{a\in A}(-1)^a=(-1)^{\sum_{a\in A}a}\prod_{a\in A}((-1)^aa)\\
&=(-1)^{\frac{p^2-1}{8}}\prod_{a\in A}((-1)^aa)
\end{align*}
ここで等差数列の和の公式から
$\sum_{a\in A}=\frac{p^2-1}{8}$
であることを用いている．
\[
B=\{2i\mid 0<2i\in A, i\in \zz\}, 
\]
でなおかつ
 $C=A\setminus B$とおく．
 $B$は$A$に属する偶数全体の集合で
 $C$は$A$
 に属する奇数の集合である．
 もちろん
 $A=B\sqcup C$である．
$D=\{2i\mid \frac{p}{2}<2i<p, i\in \zz\}$
とおく．
$D$は$\frac{p}{2}$より大きく
$p$
よりも小さい偶数全体の集合である．
簡単にわかることだが，
任意の
$d\in D$
に対して，
$c\in C$
がただ一つ存在して
$p-c=d$となる．
よって
$D$と
$C$の濃度は等しいことに注意しよう．
よって，
\[
\prod_{a\in B}(-a)=_p \prod_{2i\in D}2i=2^{\card{C}}\prod_{\frac{p}{4}<i<\frac{p}{2}}i
\]
となる．
また，
\[
\prod_{a\in B}a=\prod_{2i\in B}2i=2^{\card{B}}\prod_{0<i<\frac{p}{4}}i
\]
となる．よって，
\begin{align*}
\prod_{a\in A}((-1)^aa)&=_p\left(2^{\card{C}}\prod_{\frac{p}{4}<i<\frac{p}{2}}i\right)
\left(2^{\card{B}}\prod_{0<i<\frac{p}{4}}i\right)\\
&=_p 2^{\card(B)+\card(C)}\prod_{0<i<\frac{p}{2}}i\\
&=_p 2^{\card(A)}\prod_{a\in A}a=_p2^{\frac{p-1}{2}}\prod_{a\in A}a
\end{align*}
よって
以上の式を合わせて，
\[
\prod_{a\in A}a=_p(-1)^{\frac{p^2-1}{8}}2^{\frac{p-1}{2}}\prod_{a\in A}a
\]
を得る．
$\prod_{a\in A}a$
は
$\ff_p$の中で$0$
ではないので，
\[
2^{\frac{p-1}{2}}=_p(-1)^{\frac{p^2-1}{8}}
\]
よってオイラーの基準から，
\[
\leg{2}{p}=_p(-1)^{\frac{p^2-1}{8}}
\]
を得るが，
両辺はそれぞれ$1$か$-1$の値しか取らず，
$2\neq_p0$なので，
結局
\[
\leg{2}{p}=(-1)^{\frac{p^2-1}{8}}
\]
を得る．
これが示すべき等式であった．
\end{proof}

平方剰余の相互法則は次のものである．
\begin{prop}
2つの異なる奇素数$p,q$に対して，
\[
\leg{p}{q}\leg{q}{p}=(-1)^{\frac{p-1}{2}\frac{q-1}{2}}
\]
\end{prop}
\begin{proof}
中国剰余定理を用いた証明を紹介する．
$p, q$は異なる奇素数なので，
特に互いに素である．
よって中国剰余定理から可逆元の間の同型
\[
(\zz/pq\zz)^{\times}\cong (\zz/p\zz)^{\times}\times (\zz/q\zz)^{\times}
\]
を得る．
さらに詳しくいうと，
写像
\[
\delta：（\zz/pq\zz)^{\times}\to (\zz/p\zz)^{\times}\times (\zz/q\zz)^{\times}; [k]_{pq}\mapsto ([k]_p, [k]_q)
\]
はwell-defined だし，同型写像であるということである．
さて，
\[
A=\{[a]_{pq}\in (\zz/pq\zz)\mid 0<a<pq/2\}, 
\]
でなおかつ
\[
B=\{([a]_p, [b]_q)\in  (\zz/p\zz)^{\times}\times (\zz/q\zz)^{\times} \mid a\in \zz, 0<b<q/2\}
\]
とする．
任意の
$(a, b)\in  B$
に対して
$k\in A$
が存在して
$\delta(k)=(a, b)$
かもしくは
$\delta(k)=-(a, b)$
のどちらか一方のみが成り立つ．
よって，
ある$\epsilon\in \{1, -1\}$
が存在して．
\[
\prod_{(x, y)\in B}(x, y)=\epsilon\prod_{k\in A}\delta(k)
\]
が成り立つ．
さて，
\begin{align*}
\prod_{(x, y)\in B}(x, y)&=\prod_{0<a<p, 0<b<q/2}([a]_p, [b]_q)
=\prod_{0<b<q/2}\prod_{0<a<p}([a]_p, [b]_q)
\\
&=\prod_{0<b<q/2}([(p-1)!]_p, [b^{p-1}]_q)
=\left(\left[((p-1)!)^{\frac{q-1}{2}}\right]_p, \left[\left((\frac{q-1}{2})!\right)^{p-1}\right]_q\right)
\end{align*}
である．
ここでウィルソンの定理から$[(p-1)!]_p=[-1]_p$であり，
\[
\left(\left(\frac{q-1}{2}\right)!\right)^2=_q(q-1)!(-1)^{\frac{q-1}{2}}=_q(-1)(-1)^{\frac{q-1}{2}}
\]
なので，
\begin{equation}\label{eq:daiji}
\prod_{(x, y)\in B}(x, y)=
\left(\left[(-1)^{\frac{q-1}{2}}\right]_p, 
\left[(-1)^{\frac{p-1}{2}}(-1)^{\frac{p-1}{2}\frac{q-1}{2}}\right]_q\right)
\end{equation}
である．
そして
\begin{align*}
\prod_{k\in A}k&=\prod_{1\le k<pq/2, gcd(k, pq)=1}k\\
&=\left(\prod_{1\le k<pq/2, p\not| k}k\right)\cdot 
\left(\prod_{1\le k<pq/2,\ q|k}k\right)^{-1}\\
&=\left(\prod_{0<k<p}k\right)\left(\prod_{p<k<2q}\right)\cdots \left(\prod_{(Q-1)p<k<Qp}k\right)\left(\prod_{Qp<k<pq/2}k\right)\left(\prod_{1\le k<pq/2, q|k}k\right)^{-1}
\end{align*}
今からこの量の$p$を法とした値を考える．
\[
\prod_{(i-1)p<k<ip}k=_p (p-1)!=_p-1
\]
であり，
$\{k\in \zz\mid Qp<k<pq/2 \}=\{k\in \zz\mid Qp+1\le k\le (pq-1)/2\}=\{Qp+1, Qp+2, \dots, Qp+P\}$
なので
\[
\prod_{Qp<k<pq/2}k=_pP!
\]
また，
$\{k\in \zz\mid 1\le k<pq/2\}=\{k\in \zz\mid 1\le k\le (pq-1)/2\}
=\{1,2,\dots, (pq-1)/2\}$
であるから，
$\{k\in \zz\mid 1\le k<pq/2\}=\{q, 2q, \dots, Pq\}$
なので
\[
\prod_{1\le k<pq/2, q|k}k=\prod_{1\le i\le P}(qi)=q^{P}P!
\]
となる．
よって，
\[
\prod_{k\in A}k=_p\frac{(-1)^Q\cdot P!}{q^{P}P!}
\]
であるが，
オイラーの基準から
\[
\prod_{k\in A}k=_p(-1)^Q\leg{q}{p}
\]

今の議論を，
$p$と
$q$
の役割を入れ替えて適用すると，
同様にして
\[
\prod_{k\in A}k=_q(-1)^P\leg{p}{q}
\]
よって(\ref{eq:daiji})と比較すると，
\[
(-1)^Q\leg{q}{p}=_p\epsilon(-1)^{Q}
\]
と
\[
(-1)^P\leg{p}{q}=_q\epsilon(-1)^P(-1)^{PQ}
\]
これら二式の両辺のそれぞれの量も
$1$
か
$-1$
のどちらかの値しか取らず，
$2\neq_p0$でなおかつ
 $2\neq_q0$なので，
 実際の等式として
 \[
 \leg{q}{p}=\epsilon
 \]
 と
 \[
 \leg{p}{q}=(-1)^{PQ}
 \]
 を得る．
 この二式から，
 \[
 \leg{p}{q}\leg{q}{p}=(-1)^{\frac{p-1}{2}\frac{q-1}{2}}
 \]
 を得る．
 これが示したい等式であった．
\end{proof}





	
\begin{thebibliography}{9}
\bibitem{K}
S. Y. Kim, \textit{An elementary proof of the quadratic reciprocity law}, 
Amer. Math. Soc. 111 (2004), no. 1, 48--50. 
\bibitem{R} 
G. Rousseau, \textit{On the quadratic reciprocity law}, 
J. Austral. Math. Soc. (Series A), 51 (1991), 423--425. 
\bibitem{Ex}
Stack Exchangeの質問「Legendre symbol. second supplementary law」に対する回答, 
\begin{verbatim}https://math.stackexchange.com/questions/
180002/legendre-symbol-second-supplementary-law/180022#180022
\end{verbatim}
\end{thebibliography}




			\end{document}



\begin{comment}
[集合のやつは$\mid$を使う。]
[真なる包含関係]
\subsetneqq
[可換図式]
\[\xymatrix{
&   & (2) \ar@{=>}[rd]&  &  &  & (5) \ar@{=>}[rd]&\\ 
& (1)\ar@{=>}[ru] \ar@{=>}[rd]&  &(4)\ar@{=>}[r]&(7)& (1)\ar@{=>}[ru] \ar@{=>}[rd]& & (7)&\\ 
&   & (3) \ar@{=>}[ur]&  &  &  &(6) \ar@{=>}[ur]&
}\]

[\bigin \endを使わない方法]

{\prop[aa] aaa}
で
\begin{prop}[aa]
aaa
\end{prop}
と同じ\df \thm \proof でも同様

[直和]
$\amalg$

[同値関係でよく使う波線]
\sim

[引数付きの新命令]
\newcommand{\prn}[1]{\|#1\|_p}

[番号のリセット]

\setcounter{equation}{0}


[字体の変更]

{\rm hogehoge}と使う。

\rm ローマン
\it イタリック
\sl 斜体

[supp]

\newcommand{\supp}{\text{{\rm supp}}}

[中央揃えの改行]

	\begin{eqnarray*}
		hoge1&=&hogehoge
		hoge2&=&hogehofe

	\end{eqnarray*}

&&の部分で揃えられる。


[\tag]
\begin{equation}
f(x) = ax^2 + bx + c \tag{1.2.3}
\end{equation}



[場合分け]

	\begin{cases}
		hoge1 & \text{状況１}\\
		hoge2 & \text{状況２}\\
		・
		・
		・
	\end{cases}



\end{comment}





